%%======================================================================
%%  Chapitre 1 : Étude de l'organisme
%%======================================================================
\chapter{Étude de l'organisme}

\section{Présentation de l'organisme}

% ──────────────────────────────────────────────────────────────────────
% CONSIGNE : Remplacez le texte ci-dessous par les informations réelles
% de votre organisme d'accueil.
% ──────────────────────────────────────────────────────────────────────

L'organisme d'accueil \textbf{[Nom de l'organisme]} est un établissement spécialisé dans [domaine d'activité]. Fondé en [année de création], il est situé à [adresse / ville].

\vspace{0.5cm}

\textbf{Missions principales :}
\begin{itemize}[noitemsep]
    \item [Mission 1 de l'organisme] ;
    \item [Mission 2 de l'organisme] ;
    \item [Mission 3 de l'organisme].
\end{itemize}

\vspace{0.5cm}

\textbf{Structure organisationnelle :}

L'organisme est structuré en plusieurs départements / services :
\begin{itemize}[noitemsep]
    \item \textbf{Direction générale} : supervision et stratégie globale ;
    \item \textbf{Département informatique} : développement de solutions logicielles et maintenance des systèmes d'information ;
    \item \textbf{[Autres départements]} : [description].
\end{itemize}

% Décommentez et adaptez si vous avez un organigramme :
% \begin{figure}[H]
%     \centering
%     \includegraphics[width=0.85\textwidth]{images/organigramme.png}
%     \caption{Organigramme de l'organisme d'accueil}
%     \label{fig:organigramme}
% \end{figure}

%%======================================================================
\section{Cadre du stage}

\subsection*{Contexte}

Dans le cadre de notre formation en [filière / spécialité] à [nom de l'université], nous avons effectué un stage de fin d'études d'une durée de [durée] au sein de \textbf{[Nom de l'organisme]}.

\vspace{0.5cm}

Ce stage s'est déroulé du \textbf{[date de début]} au \textbf{[date de fin]}, au sein du département informatique de l'organisme.

\subsection*{Problématique}

L'organisme était confronté à plusieurs défis dans la gestion de ses activités pédagogiques :
\begin{itemize}[noitemsep]
    \item L'absence d'un système centralisé pour la gestion des cours, devoirs et évaluations ;
    \item La difficulté de suivi des absences et des justifications ;
    \item Le traitement manuel et lent des demandes administratives et des réclamations ;
    \item Le manque de visibilité en temps réel sur les notes et les soumissions ;
    \item L'absence de système de notification automatisé pour les délais et les annonces.
\end{itemize}

\subsection*{Objectifs du stage}

Pour répondre à cette problématique, les objectifs suivants ont été fixés :
\begin{enumerate}[noitemsep]
    \item Concevoir et développer une \textbf{plateforme web éducative} complète et moderne ;
    \item Implémenter un \textbf{backend robuste} avec une API REST exposant tous les services nécessaires ;
    \item Développer un \textbf{frontend réactif} avec des espaces dédiés pour chaque rôle (Admin, Enseignant, Étudiant) ;
    \item Mettre en place un système de \textbf{gestion des absences} avec justifications ;
    \item Développer un système de \textbf{soumission et évaluation} de devoirs ;
    \item Implémenter un système de \textbf{notifications} et de \textbf{réclamations} ;
    \item Assurer la \textbf{gestion des demandes administratives} et des demandes de stage.
\end{enumerate}

\subsection*{Méthodologie de travail}

Pour la réalisation de ce projet, nous avons adopté une approche itérative et incrémentale :
\begin{enumerate}[noitemsep]
    \item \textbf{Phase d'analyse} : étude des besoins, identification des acteurs et des cas d'utilisation ;
    \item \textbf{Phase de conception} : modélisation UML (diagrammes de classes, de séquence, de cas d'utilisation) ;
    \item \textbf{Phase de développement} : implémentation du backend (Spring Boot), puis du frontend (React/TypeScript) ;
    \item \textbf{Phase de test} : tests fonctionnels avec Postman et validation des interfaces.
\end{enumerate}

\newpage
